\IfFileExists{stacks-project.cls}{%
\documentclass{stacks-project}
}{%
\documentclass{amsart}
}

% For dealing with references we use the comment environment
\usepackage{verbatim}
\newenvironment{reference}{\comment}{\endcomment}
%\newenvironment{reference}{}{}
\newenvironment{slogan}{\comment}{\endcomment}
\newenvironment{history}{\comment}{\endcomment}

% For commutative diagrams we use Xy-pic
\usepackage[all]{xy}

% We use 2cell for 2-commutative diagrams.
\xyoption{2cell}
\UseAllTwocells

% We use multicol for the list of chapters between chapters
\usepackage{multicol}

% This is generally recommended for better output
\usepackage{lmodern}
\usepackage[T1]{fontenc}

% For cross-file-references
\usepackage{xr-hyper}

% Package for hypertext links:
\usepackage{hyperref}

% For any local file, say "hello.tex" you want to link to please
% use \externaldocument[hello-]{hello}
\externaldocument[introduction-]{introduction}
\externaldocument[conventions-]{conventions}
\externaldocument[sets-]{sets}
\externaldocument[categories-]{categories}
\externaldocument[topology-]{topology}
\externaldocument[sheaves-]{sheaves}
\externaldocument[sites-]{sites}
\externaldocument[stacks-]{stacks}
\externaldocument[fields-]{fields}
\externaldocument[algebra-]{algebra}
\externaldocument[brauer-]{brauer}
\externaldocument[homology-]{homology}
\externaldocument[derived-]{derived}
\externaldocument[simplicial-]{simplicial}
\externaldocument[more-algebra-]{more-algebra}
\externaldocument[smoothing-]{smoothing}
\externaldocument[modules-]{modules}
\externaldocument[sites-modules-]{sites-modules}
\externaldocument[injectives-]{injectives}
\externaldocument[cohomology-]{cohomology}
\externaldocument[sites-cohomology-]{sites-cohomology}
\externaldocument[dga-]{dga}
\externaldocument[dpa-]{dpa}
\externaldocument[sdga-]{sdga}
\externaldocument[hypercovering-]{hypercovering}
\externaldocument[schemes-]{schemes}
\externaldocument[constructions-]{constructions}
\externaldocument[properties-]{properties}
\externaldocument[morphisms-]{morphisms}
\externaldocument[coherent-]{coherent}
\externaldocument[divisors-]{divisors}
\externaldocument[limits-]{limits}
\externaldocument[varieties-]{varieties}
\externaldocument[topologies-]{topologies}
\externaldocument[descent-]{descent}
\externaldocument[perfect-]{perfect}
\externaldocument[more-morphisms-]{more-morphisms}
\externaldocument[flat-]{flat}
\externaldocument[groupoids-]{groupoids}
\externaldocument[more-groupoids-]{more-groupoids}
\externaldocument[etale-]{etale}
\externaldocument[chow-]{chow}
\externaldocument[intersection-]{intersection}
\externaldocument[pic-]{pic}
\externaldocument[weil-]{weil}
\externaldocument[adequate-]{adequate}
\externaldocument[dualizing-]{dualizing}
\externaldocument[duality-]{duality}
\externaldocument[discriminant-]{discriminant}
\externaldocument[derham-]{derham}
\externaldocument[local-cohomology-]{local-cohomology}
\externaldocument[algebraization-]{algebraization}
\externaldocument[curves-]{curves}
\externaldocument[resolve-]{resolve}
\externaldocument[models-]{models}
\externaldocument[functors-]{functors}
\externaldocument[equiv-]{equiv}
\externaldocument[pione-]{pione}
\externaldocument[etale-cohomology-]{etale-cohomology}
\externaldocument[proetale-]{proetale}
\externaldocument[relative-cycles-]{relative-cycles}
\externaldocument[more-etale-]{more-etale}
\externaldocument[trace-]{trace}
\externaldocument[crystalline-]{crystalline}
\externaldocument[spaces-]{spaces}
\externaldocument[spaces-properties-]{spaces-properties}
\externaldocument[spaces-morphisms-]{spaces-morphisms}
\externaldocument[decent-spaces-]{decent-spaces}
\externaldocument[spaces-cohomology-]{spaces-cohomology}
\externaldocument[spaces-limits-]{spaces-limits}
\externaldocument[spaces-divisors-]{spaces-divisors}
\externaldocument[spaces-over-fields-]{spaces-over-fields}
\externaldocument[spaces-topologies-]{spaces-topologies}
\externaldocument[spaces-descent-]{spaces-descent}
\externaldocument[spaces-perfect-]{spaces-perfect}
\externaldocument[spaces-more-morphisms-]{spaces-more-morphisms}
\externaldocument[spaces-flat-]{spaces-flat}
\externaldocument[spaces-groupoids-]{spaces-groupoids}
\externaldocument[spaces-more-groupoids-]{spaces-more-groupoids}
\externaldocument[bootstrap-]{bootstrap}
\externaldocument[spaces-pushouts-]{spaces-pushouts}
\externaldocument[spaces-chow-]{spaces-chow}
\externaldocument[groupoids-quotients-]{groupoids-quotients}
\externaldocument[spaces-more-cohomology-]{spaces-more-cohomology}
\externaldocument[spaces-simplicial-]{spaces-simplicial}
\externaldocument[spaces-duality-]{spaces-duality}
\externaldocument[formal-spaces-]{formal-spaces}
\externaldocument[restricted-]{restricted}
\externaldocument[spaces-resolve-]{spaces-resolve}
\externaldocument[formal-defos-]{formal-defos}
\externaldocument[defos-]{defos}
\externaldocument[cotangent-]{cotangent}
\externaldocument[examples-defos-]{examples-defos}
\externaldocument[algebraic-]{algebraic}
\externaldocument[examples-stacks-]{examples-stacks}
\externaldocument[stacks-sheaves-]{stacks-sheaves}
\externaldocument[criteria-]{criteria}
\externaldocument[artin-]{artin}
\externaldocument[quot-]{quot}
\externaldocument[stacks-properties-]{stacks-properties}
\externaldocument[stacks-morphisms-]{stacks-morphisms}
\externaldocument[stacks-limits-]{stacks-limits}
\externaldocument[stacks-cohomology-]{stacks-cohomology}
\externaldocument[stacks-perfect-]{stacks-perfect}
\externaldocument[stacks-introduction-]{stacks-introduction}
\externaldocument[stacks-more-morphisms-]{stacks-more-morphisms}
\externaldocument[stacks-geometry-]{stacks-geometry}
\externaldocument[moduli-]{moduli}
\externaldocument[moduli-curves-]{moduli-curves}
\externaldocument[examples-]{examples}
\externaldocument[exercises-]{exercises}
\externaldocument[guide-]{guide}
\externaldocument[desirables-]{desirables}
\externaldocument[coding-]{coding}
\externaldocument[obsolete-]{obsolete}
\externaldocument[fdl-]{fdl}
\externaldocument[index-]{index}

% Theorem environments.
%
\theoremstyle{plain}
\newtheorem{theorem}[subsection]{Theorem}
\newtheorem{proposition}[subsection]{Proposition}
\newtheorem{lemma}[subsection]{Lemma}

\theoremstyle{definition}
\newtheorem{definition}[subsection]{Definition}
\newtheorem{example}[subsection]{Example}
\newtheorem{exercise}[subsection]{Exercise}
\newtheorem{situation}[subsection]{Situation}

\theoremstyle{remark}
\newtheorem{remark}[subsection]{Remark}
\newtheorem{remarks}[subsection]{Remarks}

\numberwithin{equation}{subsection}

% Macros
%
\def\lim{\mathop{\mathrm{lim}}\nolimits}
\def\colim{\mathop{\mathrm{colim}}\nolimits}
\def\Spec{\mathop{\mathrm{Spec}}}
\def\Hom{\mathop{\mathrm{Hom}}\nolimits}
\def\Ext{\mathop{\mathrm{Ext}}\nolimits}
\def\SheafHom{\mathop{\mathcal{H}\!\mathit{om}}\nolimits}
\def\SheafExt{\mathop{\mathcal{E}\!\mathit{xt}}\nolimits}
\def\Sch{\mathit{Sch}}
\def\Mor{\mathop{\mathrm{Mor}}\nolimits}
\def\Ob{\mathop{\mathrm{Ob}}\nolimits}
\def\Sh{\mathop{\mathit{Sh}}\nolimits}
\def\NL{\mathop{N\!L}\nolimits}
\def\CH{\mathop{\mathrm{CH}}\nolimits}
\def\proetale{{pro\text{-}\acute{e}tale}}
\def\etale{{\acute{e}tale}}
\def\QCoh{\mathit{QCoh}}
\def\Ker{\mathop{\mathrm{Ker}}}
\def\Im{\mathop{\mathrm{Im}}}
\def\Coker{\mathop{\mathrm{Coker}}}
\def\Coim{\mathop{\mathrm{Coim}}}

% Boxtimes
%
\DeclareMathSymbol{\boxtimes}{\mathbin}{AMSa}{"02}

%
% Macros for moduli stacks/spaces
%
\def\QCohstack{\mathcal{QC}\!\mathit{oh}}
\def\Cohstack{\mathcal{C}\!\mathit{oh}}
\def\Spacesstack{\mathcal{S}\!\mathit{paces}}
\def\Quotfunctor{\mathrm{Quot}}
\def\Hilbfunctor{\mathrm{Hilb}}
\def\Curvesstack{\mathcal{C}\!\mathit{urves}}
\def\Polarizedstack{\mathcal{P}\!\mathit{olarized}}
\def\Complexesstack{\mathcal{C}\!\mathit{omplexes}}
% \Pic is the operator that assigns to X its picard group, usage \Pic(X)
% \Picardstack_{X/B} denotes the Picard stack of X over B
% \Picardfunctor_{X/B} denotes the Picard functor of X over B
\def\Pic{\mathop{\mathrm{Pic}}\nolimits}
\def\Picardstack{\mathcal{P}\!\mathit{ic}}
\def\Picardfunctor{\mathrm{Pic}}
\def\Deformationcategory{\mathcal{D}\!\mathit{ef}}
\usepackage{cleveref}
\crefname{originalreference}{original reference}{original references}
\Crefname{originalreference}{Original reference}{Original references}
\makeatletter
\newlabel{lemma-augmentation-site}{{0D70}{source}{}{sourcepointer-1}{}}
\newlabel{lemma-augmentation-site@cref}{{[originalreference][0][]0D70}{[1][0][]source}{}{}{}}
\makeatother
\crefname{originalreference}{original reference}{original references}
\Crefname{originalreference}{Original reference}{Original references}
\makeatletter
\newlabel{lemma-restriction-to-components-site}{{09WG}{source}{}{sourcepointer-2}{}}
\newlabel{lemma-restriction-to-components-site@cref}{{[originalreference][0][]09WG}{[1][0][]source}{}{}{}}
\makeatother
\crefname{originalreference}{original reference}{original references}
\Crefname{originalreference}{Original reference}{Original references}
\makeatletter
\newlabel{lemma-simplicial-resolution-Z-site}{{09WI}{source}{}{sourcepointer-3}{}}
\newlabel{lemma-simplicial-resolution-Z-site@cref}{{[originalreference][0][]09WI}{[1][0][]source}{}{}{}}
\makeatother
\crefname{originalreference}{original reference}{original references}
\Crefname{originalreference}{Original reference}{Original references}
\makeatletter
\newlabel{remark-augmentation-site}{{0D6Z}{source}{}{sourcepointer-4}{}}
\newlabel{remark-augmentation-site@cref}{{[originalreference][0][]0D6Z}{[1][0][]source}{}{}{}}
\makeatother
\crefname{originalreference}{original reference}{original references}
\Crefname{originalreference}{Original reference}{Original references}
\makeatletter
\newlabel{situation-simplicial-site}{{09WE}{source}{}{sourcepointer-5}{}}
\newlabel{situation-simplicial-site@cref}{{[originalreference][0][]09WE}{[1][0][]source}{}{}{}}
\makeatother
\title{The simplicial resolution of a\textasciicircum{}\{-1\}G by the shriek functors of the augmentation}
\author{The Stacks Project Authors}
\date{stacks-project:spaces-simplicial.tex:local-snapshot-20261012; GNU FDL 1.2 or later}
\begin{document}
\maketitle
\noindent\textit{Editorial note.} All delivered mathematical text is the author's own, copied from the pinned Stacks Project chapter file sources/stacks-spaces-simplicial/pinned-original-source.tex; no statement or proof step was rewritten, reordered or removed.\par\medskip
\noindent\textit{Editorial note.} The counted proof is the author's complete argument: the opening 'We encourage the reader to read the proof of Lemma simplicial-resolution-Z-site first' (source lines 1945-1947) is a reading suggestion, not an omission - the argument that follows is complete and its conclusion rests on the homotopy equivalence proven in that lemma; the description of the functors g\_\{n!\} used 'without further mention' is established by card L46149, and the augmentation machinery by card L46155 of this batch.\par\medskip
\noindent\textit{Editorial note.} The label lemma-simplicial-resolution-augmentation names the counted statement, and the proof environment immediately following it in the source is delivered verbatim as its proof.\par\medskip
\noindent\textit{Editorial note.} The counted statement is carried by the author's own lemma environment in the chapter source.\par\medskip
\noindent\textit{Allowed standard prerequisites (editorial).}\begin{itemize}
\item The augmentation datum and its topos morphisms. Same chapter, \textbackslash{}label\{lemma-augmentation-site\} (card L46155 of this batch) and Remark \textbackslash{}label\{remark-augmentation-site\}. Lines 1936-1937: the morphisms a\_n, a and the adjunctions a\_n\textasciicircum{}\{-1\} -| a\_\{n,*\} used to define the maps.
\item The component shriek functors and their formulas. Same chapter, \textbackslash{}label\{lemma-restriction-to-components-site\} (card L46149 of this batch). Lines 1937-1943: g\_\{n!\}(a\_n\textasciicircum{}\{-1\}G) restricted to C\_m is the direct sum over varphi : [n] -> [m] of f\_varphi\textasciicircum{}\{-1\}a\_n\textasciicircum{}\{-1\}G = a\_m\textasciicircum{}\{-1\}G.
\item The simplicial resolution of the constant sheaf Z. Same chapter, \textbackslash{}label\{lemma-simplicial-resolution-Z-site\}. Lines 1966-1985: the g\_m\textasciicircum{}\{-1\}-image of the complex is a\_m\textasciicircum{}\{-1\}G tensor the Z-complex, whose homotopy equivalence with Z in degree 0 was proven there.
\item The cosimplicial identities of the coface maps. Simplicial chapter, the cosimplicial identities of the delta\textasciicircum{}n\_i. Lines 1946-1952: the alternating-sum maps d\textasciicircum{}n\_i are adjoints to the coproduct insertions labeled by delta\textasciicircum{}n\_i.
\end{itemize}\medskip
\ifdefined\hypertarget\hypertarget{sourcepointer-1}{}\fi
\noindent\textit{Source reference.} Original reference 0D70: sources/stacks-spaces-simplicial/pinned-original-source.tex, source lines 1141--1141. Prerequisite pointer invoked at proof line(s) 1948; it names a standard fact of the same chapter that the counted argument uses, and every step of the counted proof built on it is retained verbatim in the delivered body.\par\medskip
\ifdefined\hypertarget\hypertarget{sourcepointer-2}{}\fi
\noindent\textit{Source reference.} Original reference 09WG: sources/stacks-spaces-simplicial/pinned-original-source.tex, source lines 970--970. Prerequisite pointer invoked at proof line(s) 1950; it names a standard fact of the same chapter that the counted argument uses, and every step of the counted proof built on it is retained verbatim in the delivered body.\par\medskip
\ifdefined\hypertarget\hypertarget{sourcepointer-3}{}\fi
\noindent\textit{Source reference.} Original reference 09WI: sources/stacks-spaces-simplicial/pinned-original-source.tex, source lines 1711--1711. Prerequisite pointer invoked at proof line(s) 1947; it names a standard fact of the same chapter that the counted argument uses, and every step of the counted proof built on it is retained verbatim in the delivered body.\par\medskip
\ifdefined\hypertarget\hypertarget{sourcepointer-4}{}\fi
\noindent\textit{Source reference.} Original reference 0D6Z: sources/stacks-spaces-simplicial/pinned-original-source.tex, source lines 1125--1125. Prerequisite pointer invoked at proof line(s) 1932; it names a standard fact of the same chapter that the counted argument uses, and every step of the counted proof built on it is retained verbatim in the delivered body.\par\medskip
\ifdefined\hypertarget\hypertarget{sourcepointer-5}{}\fi
\noindent\textit{Source reference.} Original reference 09WE: sources/stacks-spaces-simplicial/pinned-original-source.tex, source lines 888--888. Prerequisite pointer invoked at proof line(s) 1930; it names a standard fact of the same chapter that the counted argument uses, and every step of the counted proof built on it is retained verbatim in the delivered body.\par\medskip
\begin{lemma}
\label{lemma-simplicial-resolution-augmentation}
In Situation \ref{situation-simplicial-site} let
$a_0$ be an augmentation towards a site $\mathcal{D}$
as in Remark \ref{remark-augmentation-site}.
For any abelian sheaf $\mathcal{G}$ on $\mathcal{D}$ 
there is an exact complex
$$
\ldots \to
g_{2!}(a_2^{-1}\mathcal{G}) \to
g_{1!}(a_1^{-1}\mathcal{G}) \to
g_{0!}(a_0^{-1}\mathcal{G}) \to
a^{-1}\mathcal{G} \to 0
$$
of abelian sheaves on $\mathcal{C}_{total}$.
\end{lemma}

\begin{proof}
We encourage the reader to read the proof of
Lemma \ref{lemma-simplicial-resolution-Z-site} first.
We will use Lemma \ref{lemma-augmentation-site} and
the description of the functors $g_{n!}$ in
Lemma \ref{lemma-restriction-to-components-site} without further mention.
In particular $g_{n!}(a_n^{-1}\mathcal{G})$ is the
sheaf on $\mathcal{C}_{total}$ whose restriction to $\mathcal{C}_m$
is the sheaf
$$
\bigoplus\nolimits_{\varphi : [n] \to [m]} f_\varphi^{-1}a_n^{-1}\mathcal{G} =
\bigoplus\nolimits_{\varphi : [n] \to [m]} a_m^{-1}\mathcal{G}
$$
As maps of the complex we take $\sum (-1)^i d^n_i$ where
$d^n_i : g_{n!}(a_n^{-1}\mathcal{G}) \to g_{n - 1!}(a_{n - 1}^{-1}\mathcal{G})$
is the adjoint to the map
$a_n^{-1}\mathcal{G} \to \bigoplus_{[n - 1] \to [n]} a_n^{-1}\mathcal{G} =
g_n^{-1}g_{n - 1!}(a_{n - 1}^{-1}\mathcal{G})$
corresponding to the factor labeled with $\delta^n_i : [n - 1] \to [n]$.
The map $g_{0!}(a_0^{-1}\mathcal{G}) \to a^{-1}\mathcal{G}$ is adjoint
to the identity map of $a_0^{-1}\mathcal{G}$.
Then $g_m^{-1}$ applied to the chain complex in degrees
$\ldots, 2, 1, 0$ gives the complex
$$
\ldots \to
\bigoplus\nolimits_{\alpha \in \Mor_\Delta([2], [m])]} a_m^{-1}\mathcal{G} \to
\bigoplus\nolimits_{\alpha \in \Mor_\Delta([1], [m])]} a_m^{-1}\mathcal{G} \to
\bigoplus\nolimits_{\alpha \in \Mor_\Delta([0], [m])]} a_m^{-1}\mathcal{G}
$$
on $\mathcal{C}_m$. This is equal to $a_m^{-1}\mathcal{G}$
tensored over the constant sheaf $\mathbf{Z}$ with the complex
$$
\ldots \to
\bigoplus\nolimits_{\alpha \in \Mor_\Delta([2], [m])]} \mathbf{Z} \to
\bigoplus\nolimits_{\alpha \in \Mor_\Delta([1], [m])]} \mathbf{Z} \to
\bigoplus\nolimits_{\alpha \in \Mor_\Delta([0], [m])]} \mathbf{Z}
$$
discussed in the proof of Lemma \ref{lemma-simplicial-resolution-Z-site}.
There we have seen that this complex is homotopy equivalent to
$\mathbf{Z}$ placed in degree $0$ which finishes the proof.
\end{proof}

\end{document}
